
\documentclass[10pt,a4paper]{article}

% ============================================================
% PACKAGES
% ============================================================
\usepackage[T1]{fontenc}
\usepackage[utf8]{inputenc}
\usepackage[french]{babel}
\usepackage{lmodern}
\usepackage{microtype}

%\usepackage[a4paper,margin=1.35cm,headheight=13pt]{geometry}
\usepackage{fancyhdr}
\usepackage{titlesec}
\usepackage{enumitem}
\usepackage{booktabs}
\usepackage{array}
\usepackage{amsmath}
\usepackage{xcolor}
\usepackage{caption}
\usepackage{tikz}

\usetikzlibrary{
    positioning,
    arrows.meta,
    fit,
    backgrounds,
    calc,
    decorations.pathreplacing,
    shapes.geometric
}

\usepackage[
    a4paper,
    left=1.7cm,
    right=1.7cm,
    top=1.7cm,
    bottom=1.6cm,
    headheight=14pt,
    headsep=7pt,
    footskip=12pt
]{geometry}

% ============================================================
% COULEURS
% ============================================================
\definecolor{encre}{HTML}{16222E}
\definecolor{accent}{HTML}{1F6F8B}
\definecolor{accentpale}{HTML}{E7F0F4}
\definecolor{gris}{HTML}{6E7C87}
\definecolor{vert}{HTML}{2E7D5B}
\definecolor{vertpale}{HTML}{E6F2EC}

\usepackage[colorlinks=true,linkcolor=accent,urlcolor=accent]{hyperref}

% ============================================================
% STYLE
% ============================================================
\titleformat{\section}
    {\normalfont\large\bfseries\sffamily\color{accent}}
    {\thesection}{0.7em}{}

\titleformat{\subsection}
    {\normalfont\normalsize\bfseries\sffamily\color{encre}}
    {\thesubsection}{0.6em}{}

\titlespacing*{\section}{0pt}{5pt}{2pt}
\titlespacing*{\subsection}{0pt}{3pt}{1pt}

\pagestyle{fancy}
\fancyhf{}
\fancyhead[L]{\sffamily\footnotesize\color{gris}
Déploiement d'un LLM open source}
\fancyhead[R]{\sffamily\footnotesize\color{gris}Projet DevOps}
\fancyfoot[C]{\sffamily\footnotesize\color{gris}\thepage}

\renewcommand{\headrulewidth}{0.3pt}
\renewcommand{\headrule}{
    \hbox to\headwidth{
        \color{accent!30}
        \leaders\hrule height \headrulewidth\hfill
    }
}

\setlist[itemize]{
    leftmargin=1.05em,
    itemsep=0.2pt,
    topsep=1pt,
    parsep=0pt
}

\setlength{\parindent}{0pt}
\setlength{\parskip}{2pt}

\captionsetup{
    font=small,
    labelfont={bf,sf},
    textfont=sf,
    skip=2pt
}

% ============================================================
% DOCUMENT
% ============================================================
\begin{document}

% ============================================================
% TITRE
% ============================================================
\begin{center}
{\color{accent}\rule{\textwidth}{1pt}}\\[3pt]

{\sffamily\LARGE\bfseries\color{encre}
Déploiement d'un modèle de langage open source}\\[2pt]

%{\sffamily\large\color{accent}
%De Docker compose à Kubernetes}\\[2pt]

{\sffamily\small\color{gris}
Mouhamadou Aliou BA
\quad\textbullet\quad
Projet DevOps
\quad\textbullet\quad
\today}\\[2pt]

{\color{accent}\rule{\textwidth}{1pt}}
\end{center}

% ============================================================
% 1. ARCHITECTURE
% ============================================================
\section{L'architecture déployée}
\thispagestyle{empty}
L'application réunit deux services complémentaires. Ollama héberge les modèles
de langage et expose une API d'inférence. Open WebUI fournit l'interface de
dialogue et interroge cette API. Le déploiement a d'abord été réalisé avec
Docker Compose, puis reproduit sur un cluster Minikube où neuf ressources
décrivent l'ensemble du système.

\begin{figure}[h]
\centering
\begin{tikzpicture}[
    font=\sffamily\scriptsize,
    bloc/.style={
        rectangle, rounded corners=2pt,
        draw=accent, fill=accentpale,
        minimum width=2.1cm,
        minimum height=0.62cm,
        align=center, thick,
        text=encre
    },
    pod/.style={
        rectangle, rounded corners=2pt,
        draw=encre, fill=white,
        minimum width=2.1cm,
        minimum height=0.62cm,
        align=center, thick,
        text=encre
    },
    disque/.style={
        cylinder,
        shape border rotate=90,
        aspect=0.22,
        draw=vert,
        fill=vertpale,
        minimum width=1.05cm,
        minimum height=0.85cm,
        align=center,
        text=encre
    },
    ext/.style={
        rectangle, rounded corners=2pt,
        draw=gris, fill=white,
        minimum width=1.45cm,
        minimum height=0.52cm,
        align=center
    },
    conf/.style={
        rectangle, rounded corners=2pt,
        draw=gris, fill=white,
        minimum width=1.5cm,
        minimum height=0.52cm,
        align=center,
        text=gris
    },
    fl/.style={
        -{Stealth[length=1.6mm]},
        thick,
        draw=accent
    },
    flv/.style={
        -{Stealth[length=1.6mm]},
        thick,
        draw=vert
    },
    flg/.style={
        -{Stealth[length=1.5mm]},
        thick,
        draw=gris,
        dashed
    }
]

\node[ext] (nav) {Navigateur};
\node[bloc, right=1.15cm of nav] (ing)
    {Ingress nginx\\[-1pt]\textit{llm.local}};
\node[bloc, right=1.15cm of ing] (svcw)
    {Service\\[-1pt]\textit{open-webui}};
\node[pod, below=0.55cm of svcw] (podw)
    {Pod\\[-1pt]Open WebUI};

\node[bloc, left=1.65cm of podw] (svco)
    {Service\\[-1pt]\textit{ollama}};
\node[pod, left=1.15cm of svco] (podo)
    {Pod\\[-1pt]Ollama};

\node[disque, right=0.45cm of podw] (pvcw)
    {PVC\\[-1pt]2 Gi};
\node[disque, left=0.45cm of podo] (pvco)
    {PVC\\[-1pt]8 Gi};

\node[conf, below=0.5cm of svco] (cm)
    {ConfigMap};

\draw[fl] (nav) -- (ing);
\draw[fl] (ing) -- (svcw);
\draw[fl] (svcw) -- (podw);
\draw[fl] (podw) --
    node[midway,above=1pt,font=\tiny,text=gris]{résolution DNS}
    (svco);
\draw[fl] (svco) -- (podo);
\draw[flv] (podw) -- (pvcw);
\draw[flv] (podo) -- (pvco);
\draw[flg] (cm) -- (podo);
\draw[flg] (cm) -| (podw);

\begin{scope}[on background layer]
\node[
    draw=gris,
    dashed,
    rounded corners=3pt,
    inner sep=5pt,
    fit=(ing)(svcw)(podw)(svco)(podo)(pvcw)(pvco)(cm),
    label={
        [font=\sffamily\tiny,text=gris]
        above right:namespace \texttt{llm}
    }
] {};
\end{scope}

\end{tikzpicture}
\caption{Les neuf ressources et le chemin d'une requête, du navigateur au modèle}
\end{figure}

\vspace{-3pt}

\subsection{Rôle de chaque ressource}

Le \textbf{Namespace} cloisonne les objets du projet, isole les noms et rend la
suppression complète triviale. La \textbf{ConfigMap} externalise la
configuration, injectée dans les deux conteneurs par \texttt{envFrom} plutôt
qu'inscrite dans les images ou répétée dans chaque manifeste. Les deux
\textbf{Deployments} maintiennent l'état désiré, recréant un pod dès qu'il
disparaît ou échoue. Les deux \textbf{Services} donnent une adresse stable aux
pods, dont les adresses IP changent à chaque recréation. Les deux
\textbf{PersistentVolumeClaims} réclament le stockage durable, huit gigaoctets
pour les poids des modèles et deux pour la base de l'interface. L'
\textbf{Ingress} enfin expose l'interface hors du cluster, les Services de
type ClusterIP n'étant joignables que de l'intérieur.

\subsection{Comment les services communiquent}

Trois mécanismes se superposent. Le DNS interne attribue à chaque service un nom
résoluble dans son namespace, si bien qu'Open WebUI s'adresse au nom
\texttt{ollama} sans jamais connaître l'adresse du pod. Le service relaie
ensuite le trafic vers les pods portant l'étiquette attendue et déclarés
prêts, en mettant seul ses points de terminaison à jour lorsqu'un pod est
remplacé. C'est ce découplage que Docker Compose n'offre pas au même degré.
L'Ingress assure enfin la traduction depuis l'extérieur, routant l'hôte
\texttt{llm.local} vers le service de l'interface.

\subsection{Pourquoi des volumes persistants}

Le système de fichiers d'un conteneur disparaît avec lui, or un pod est par
conception une entité éphémère que le deployment remplace à la moindre mise à
jour ou au moindre échec. Sans PVC, Ollama retéléchargerait plusieurs
gigaoctets de poids à chaque redémarrage et Open WebUI perdrait sa base, donc
ses conversations et ses réglages. Le PVC exprime un besoin que la StorageClass
du cluster satisfait, ce qui permet au même manifeste de fonctionner sur
Minikube comme chez un fournisseur cloud.

Les deux volumes sont en \texttt{ReadWriteOnce}, montables en écriture par un
seul n\oe ud. Une mise à jour progressive créerait le nouveau pod avant la
libération du volume et resterait bloquée en attente. Les deux Deployments
déclarent donc explicitement \texttt{strategy: Recreate}.

\subsection{Le rôle des trois sondes}

\begin{center}
\small
\renewcommand{\arraystretch}{1.02}
\begin{tabular}{
@{}
>{\raggedright\arraybackslash}p{2.25cm}
>{\raggedright\arraybackslash}p{4.2cm}
>{\raggedright\arraybackslash}p{8.6cm}
@{}}
\toprule
\textbf{Sonde} & \textbf{Question posée} & \textbf{Effet d'un échec} \\
\midrule
\texttt{startupProbe}
& Le démarrage est-il terminé ?
& Redémarrage du conteneur. Les deux autres sondes restent suspendues tant
qu'elle n'a pas réussi. \\

\texttt{readinessProbe}
& Peut-il recevoir du trafic ?
& Retrait des points de terminaison du Service, sans redémarrage. \\

\texttt{livenessProbe}
& Est-il encore vivant ?
& Le conteneur est tué puis redémarré. \\
\bottomrule
\end{tabular}
\end{center}

La sonde de démarrage protège les applications lentes à s'initialiser et permet
de conserver une surveillance stricte une fois le service établi. Un incident
l'a illustré avec une limite mémoire de 1 Gi et une sonde de démarrage
plafonnée à cinq minutes, le pod Open WebUI alternait \texttt{OOMKilled} et
échec de sonde sans jamais finir de démarrer, le conteneur chargeant un modèle
d'embedding au boot. Le passage à 2 Gi et dix minutes a levé les deux causes.
Une sonde de démarrage doit être dimensionnée sur le pire cas de démarrage à
froid, sans quoi elle transforme une initialisation lente en boucle d'échecs.

% ============================================================
% 2. COMPOSE VS KUBERNETES
% ============================================================
\section{Ce que Kubernetes change par rapport à Compose}

Les deux outils décrivent une application multi conteneurs de façon
déclarative, mais ne répondent pas à la même question. Compose explique comment
démarrer des conteneurs ensemble sur une machine. Kubernetes décrit un état à
maintenir quoi qu'il advienne, vers lequel des contrôleurs convergent en
permanence. Compose redémarre un conteneur défaillant là où Kubernetes recrée
un pod, éventuellement ailleurs dans le cluster. La mise à l'échelle passe
d'une option de ligne de commande à un objet déclaratif qu'un autoscaler peut
piloter. La configuration passe de variables d'environnement à des objets
versionnés, ConfigMap et Secret. Le stockage passe de volumes locaux à une
abstraction du fournisseur.

La surveillance de l'état des conteneurs, enfin, passe d'un test unique exécuté
à l'intérieur du conteneur à trois sondes pilotées par le kubelet, et cette
différence s'est vérifiée concrètement. Le contrôle fourni dans le sujet appelle
\texttt{wget}, absent de l'image Open WebUI, si bien que le conteneur restait
marqué en défaut alors qu'il répondait correctement. Les sondes
\texttt{httpGet} de Kubernetes ignorent ce problème, étant exécutées depuis
l'extérieur du conteneur sans dépendre d'un binaire embarqué.

% ============================================================
% 3. SUPERVISION
% ============================================================
\section{Superviser le déploiement}

Un déploiement conforme ne dit rien de son comportement sous charge. La partie
supervision vise donc à mesurer ce que vaut réellement le service rendu, en
termes de latence et de débit.

\subsection{La chaîne d'observabilité mise en place}

Trois outils la composent.\textbf{Prometheus} collecte périodiquement les
métriques exposées par ses cibles et les conserve en séries temporelles.\textbf{Grafana} interroge Prometheus en PromQL et rend les tableaux de bord,
la source de données et les tableaux étant provisionnés par fichier pour
rester reproductibles. Entre les deux manquait l'essentiel car, dans la version
et la configuration utilisées, \textbf{Ollama ne fournissait pas directement
les métriques Prometheus nécessaires} pour mesurer le TTFT, le TPOT et les
performances d'inférence. Il n'y avait donc rien à collecter sous la forme
attendue.

Un exporter a été écrit pour combler ce vide. Plutôt que de lire le moteur,
il le mesure de l'extérieur en lui adressant de vraies requêtes en streaming
et en les chronométrant. C'est une sonde active, dans l'esprit de
\texttt{blackbox\_exporter}. Son pool de générateurs de charge est
redimensionnable pendant l'exécution, ce qui permet d'observer la dégradation
en direct.

\begin{figure}[h]
\centering
\begin{tikzpicture}[
    font=\sffamily\scriptsize,
    b/.style={
        rectangle, rounded corners=2pt,
        draw=accent, fill=accentpale,
        minimum width=2.25cm,
        minimum height=0.62cm,
        align=center, thick,
        text=encre
    },
    e/.style={
        rectangle, rounded corners=2pt,
        draw=vert, fill=vertpale,
        minimum width=2.45cm,
        minimum height=0.62cm,
        align=center, thick,
        text=encre
    },
    fl/.style={
        -{Stealth[length=1.6mm]},
        thick,
        draw=accent
    }
]

\node[b] (ollama) {Ollama\\[-1pt]\textit{moteur}};
\node[e, right=1.35cm of ollama] (exp)
    {Exporter\\[-1pt]\textit{sonde active}};
\node[b, right=1.35cm of exp] (prom)
    {Prometheus\\[-1pt]\textit{séries}};
\node[b, right=1.35cm of prom] (graf)
    {Grafana\\[-1pt]\textit{tableaux}};

\draw[fl,draw=vert,transform canvas={yshift=3pt}]
    (exp) -- node[above,font=\tiny,text=gris]{requêtes streaming} (ollama);

\draw[fl,draw=vert,transform canvas={yshift=-3pt}]
    (ollama) -- node[below,font=\tiny,text=gris]{jetons et compteurs} (exp);

\draw[fl]
    (prom) -- node[above,font=\tiny,text=gris]{\texttt{/metrics}} (exp);

\draw[fl]
    (prom) -- node[above,font=\tiny,text=gris]{PromQL} (graf);

\end{tikzpicture}
\caption{L'exporter mesure le moteur faute de pouvoir l'interroger,
Prometheus le collecte}
\end{figure}

\vspace{-3pt}

\subsection{Les grandeurs mesurées}

\begin{figure}[h]
\centering
\begin{tikzpicture}[font=\sffamily\scriptsize, x=1cm, y=0.8cm]

\draw[
    thick,
    draw=gris,
    -{Stealth[length=1.7mm]}
]
(0,0) -- (13.3,0);

\foreach \x in {0,3.3,11.0}
{
    \draw[thick,draw=encre]
    (\x,-0.15) -- (\x,0.15);
}

\foreach \x in {4.5,5.7,6.9,8.1,9.3,10.5}
{
    \draw[draw=accent,thick]
    (\x,-0.09) -- (\x,0.09);
    \fill[accent] (\x,0.26) circle (1.2pt);
}

\fill[accent] (3.3,0.26) circle (1.2pt);

\node[below=2pt,text=encre] at (0,-0.15) {$t_{env}$};
\node[below=2pt,text=encre] at (3.3,-0.15) {$t_1$};
\node[below=2pt,text=encre] at (11.0,-0.15) {$t_{fin}$};

\draw[
    decorate,
    decoration={brace,amplitude=4pt},
    draw=accent,
    thick
]
(0,0.58) -- (3.3,0.58)
node[midway,above=3pt,text=accent]{\textbf{TTFT}};

\draw[
    decorate,
    decoration={brace,amplitude=4pt},
    draw=vert,
    thick
]
(4.5,0.58) -- (5.7,0.58)
node[midway,above=3pt,text=vert]{\textbf{TPOT}};

\draw[
    decorate,
    decoration={brace,mirror,amplitude=4pt},
    draw=gris,
    thick
]
(0,-0.72) -- (11.0,-0.72)
node[midway,below=3pt,text=gris]{\textbf{latence de bout en bout}};

\node[
    text=gris,
    anchor=west
]
at (0.1,1.42)
{\footnotesize attente, chargement du modèle, traitement du prompt};

\node[
    text=gris,
    anchor=west
]
at (6.6,1.0)
{\footnotesize génération des $N$ jetons};

\end{tikzpicture}
\caption{Décomposition temporelle d'une requête}
\end{figure}

\vspace{-3pt}

En notant $t_{env}$ l'instant d'envoi, $t_1$ l'arrivée du premier jeton,
$t_{fin}$ la fin du flux et $N$ le nombre de jetons générés, les grandeurs
suivies se définissent ainsi.

\begin{align*}
\text{TTFT}
&=
t_1-t_{env}
=
T_{attente}+T_{chargement}+T_{prefill}
\\[-1pt]
\text{TPOT}
&=
\frac{t_{fin}-t_1}{N-1}
\\[-1pt]
\text{E2E}
&=
\text{TTFT}+\text{TPOT}(N-1)
\\[-1pt]
\text{Débit}
&=
\frac{\texttt{eval\_count}}
{\texttt{eval\_duration}\times10^{-9}}
\end{align*}

Le TTFT correspond à ce que perçoit l'utilisateur avant le premier caractère.
Sa décomposition explique une observation frappante, la première requête après
démarrage affichant 123 secondes contre 0,5 seconde ensuite. Cet écart est
compatible avec un coût important de $T_{chargement}$, sans permettre à lui
seul d'isoler chaque composante du cold start.

Le PVC d'Ollama économise surtout le téléchargement répété des poids lors
d'une recréation ; il ne supprime pas le chargement des poids en mémoire lors
du redémarrage du processus.

Le débit est repris des compteurs renvoyés par le moteur dans son dernier
segment, ce qui donne une mesure serveur sans bruit réseau, tandis que le TTFT
est mesuré côté client, ce qui impose le streaming. En réponse unique, le
premier et le dernier jeton arrivent ensemble et la grandeur perd tout sens.

Prometheus ne conserve pas les valeurs individuelles mais des compteurs
cumulés par intervalle, les percentiles étant reconstruits par interpolation
linéaire dans l'intervalle où le cumul franchit le seuil, au moyen de
\texttt{histogram\_quantile}. La précision dépend donc du découpage retenu, et
les intervalles de l'exporter s'étendent de 0,25 à 120 secondes, les valeurs
par défaut regroupant toutes ces mesures dans un seul seau.

\subsection{Résultats et comportement sous concurrence}

\begin{center}
\small
\renewcommand{\arraystretch}{1.02}
\begin{tabular}{@{}lrrl@{}}
\toprule
&
\textbf{1 requête}
&
\textbf{3 en parallèle}
&
\\
\midrule
TTFT p50
& 0,5 s
& 35,0 s
& \\

TTFT p95
& --
& 57,0 s
& \\

Latence e2e p95
& 17 s
& 86,0 s
& \textcolor{accent}{$\times$ 5} \\

\textbf{Débit}
& \textbf{3,8 jetons/s}
& \textbf{3,28 jetons/s}
& \textcolor{vert}{-13,7\,\%} \\
\bottomrule
\end{tabular}
\end{center}

La latence augmente fortement alors que le débit diminue légèrement. Ces
résultats constituent des \textbf{signes de saturation dans la configuration
expérimentale}, mais les deux niveaux de concurrence testés ne permettent pas
d'identifier un point de saturation précis.

La loi de Little formalise le lien entre concurrence et temps de réponse.
En boucle fermée avec $N$ requêtes présentes dans le système, un débit de
requêtes $X$ et un temps de réponse moyen $R$, on a

\[
N=X\times R.
\]

Il serait incorrect d'utiliser directement les 3,8 et 3,28 jetons/s comme
$X$, car il s'agit d'un débit de jetons et non de requêtes. Sans débit de
requêtes mesuré, Little sert donc ici de cadre d'interprétation plutôt que
de calcul numérique de $R$.

Les mesures montrent néanmoins qu'ajouter de la concurrence augmente fortement
l'attente avant le premier jeton sans produire d'augmentation proportionnelle
du débit global. Elles révèlent ainsi une limite du \textbf{moteur d'inférence
dans la configuration étudiée}, sans permettre d'affirmer que les requêtes
sont strictement traitées en série. Améliorer le déploiement ne passe donc pas
nécessairement par un réglage Kubernetes supplémentaire, mais peut conduire
à envisager un moteur mieux adapté au service concurrent.

% ============================================================
% 4. VLLM
% ============================================================
\section{Un moteur d'inférence de production}

C'est exactement le problème que résolvent les moteurs conçus pour le service
à grande échelle.vLLM s'appuie sur deux mécanismes qui répondent chacun à
un aspect de la gestion de la concurrence.

\textbf{PagedAttention} améliore l'utilisation de la mémoire. Le cache KV,
qui mémorise les clés et valeurs d'attention des jetons déjà traités, est
habituellement alloué de façon contiguë. PagedAttention découpe le cache en
blocs de taille fixe placés librement, une table de blocs reliant chaque
position logique à son bloc physique.

Dans les expériences de l'article, cette organisation réduit fortement le
gaspillage mémoire et permet de maintenir davantage de requêtes actives en
parallèle. Des blocs identiques peuvent aussi être partagés entre requêtes
en copie sur écriture, notamment lorsque les prompts partagent un préfixe.

\textbf{Le batching continu} agit sur la file d'attente.Un lot classique est
figé, on attend la fin de toutes ses requêtes avant d'en admettre de nouvelles,
si bien que le lot avance au rythme de sa séquence la plus longue et que les
places libérées restent inutilisées.vLLM planifie à chaque itération de
décodage, une séquence terminée cédant aussitôt sa place à une requête en
attente. La concurrence peut ainsi être mieux convertie en utilisation des
ressources et en débit lorsque la charge est suffisante.

\begin{figure}[h]
\centering
\begin{tikzpicture}[
    font=\sffamily\scriptsize,
    x=1cm,
    y=0.5cm,
    req/.style={
        rectangle,
        draw=accent,
        fill=accentpale,
        minimum height=0.34cm,
        inner sep=1pt,
        text=encre
    },
    vide/.style={
        rectangle,
        draw=gris!45,
        fill=gris!12,
        minimum height=0.34cm,
        inner sep=1pt,
        text=gris
    },
    neuf/.style={
        rectangle,
        draw=vert,
        fill=vertpale,
        minimum height=0.34cm,
        inner sep=1pt,
        text=encre
    }
]

% Lot figé
\node[
    anchor=west,
    text=encre,
    font=\sffamily\scriptsize\bfseries
]
at (-0.35,3.3) {Lot figé};

\foreach \y/\len in {
    2.3/3.4,
    1.3/1.7,
    0.3/2.6,
    -0.7/1.1
}
{
    \node[
        req,
        anchor=west,
        minimum width=\len cm
    ] at (0,\y) {};
}

\foreach \y/\st/\len in {
    1.3/1.7/2.0,
    0.3/2.6/1.1,
    -0.7/1.1/2.6
}
{
    \node[
        vide,
        anchor=west,
        minimum width=\len cm
    ] at (\st,\y) {};
}

\draw[dashed,draw=gris]
(3.75,-1.2) -- (3.75,3);

\node[
    neuf,
    anchor=west,
    minimum width=1.5cm
] at (3.9,2.3) {};

\node[
    neuf,
    anchor=west,
    minimum width=1.5cm
] at (3.9,1.3) {};

% Batching continu
\begin{scope}[xshift=7.2cm]

\node[
    anchor=west,
    text=encre,
    font=\sffamily\scriptsize\bfseries
]
at (-0.35,3.3) {Batching continu};

\foreach \y/\len in {
    2.3/3.4,
    1.3/1.7,
    0.3/2.6,
    -0.7/1.1
}
{
    \node[
        req,
        anchor=west,
        minimum width=\len cm
    ] at (0,\y) {};
}

\node[
    neuf,
    anchor=west,
    minimum width=1.6cm
] at (1.7,1.3) {};

\node[
    neuf,
    anchor=west,
    minimum width=1.4cm
] at (2.6,0.3) {};

\node[
    neuf,
    anchor=west,
    minimum width=2.2cm
] at (1.1,-0.7) {};

\node[
    neuf,
    anchor=west,
    minimum width=1.1cm
] at (3.3,1.3) {};

\end{scope}

% Axes du temps
\draw[
    -{Stealth[length=1.4mm]},
    draw=gris
]
(-0.35,-1.55) -- (5.15,-1.55);

\draw[
    -{Stealth[length=1.4mm]},
    draw=gris
]
(6.85,-1.55) -- (12.35,-1.55);

\node[
    text=gris,
    font=\tiny,
    anchor=west
]
at (5.25,-1.55) {temps};

\node[
    text=gris,
    font=\tiny,
    anchor=west
]
at (12.45,-1.55) {temps};

% Légende
\node[
    req,
    minimum width=0.35cm
] (l1) at (0.35,-2.25) {};

\node[
    anchor=west,
    text=gris,
    font=\tiny
]
at (0.6,-2.25) {requêtes du lot initial};

\node[
    neuf,
    minimum width=0.35cm
] (l2) at (4.35,-2.25) {};

\node[
    anchor=west,
    text=gris,
    font=\tiny
]
at (4.6,-2.25) {requêtes admises ensuite};

\node[
    vide,
    minimum width=0.35cm
] (l3) at (9.0,-2.25) {};

\node[
    anchor=west,
    text=gris,
    font=\tiny
]
at (9.25,-2.25) {emplacement inutilisé};

\end{tikzpicture}
\caption{Quatre emplacements d'exécution. À gauche le lot attend sa séquence
la plus longue, à droite les requêtes en attente comblent les places dès
qu'elles se libèrent}
\end{figure}

\vspace{-3pt}

S'y ajoutent le cache de préfixes, qui évite de recalculer le prefill d'un
prompt système répété, et la publication native des métriques, dont la
profondeur de file d'attente, sur laquelle un autoscaling horizontal peut
directement s'appuyer.

% ============================================================
% 5. CONCLUSION
% ============================================================
\section{Conclusion}

Le déploiement a été mené à son terme dans les deux environnements, et les
ressources exigées se comportent comme prévu, sans pod en échec ni volume non
lié. L'essentiel de ce qui a été appris ne vient pourtant pas de la rédaction
des manifestes, relativement directe, mais de deux confrontations avec le réel.

La première tient au diagnostic des pannes car chacune avait une cause précise et
vérifiable, un binaire absent d'une image, une limite mémoire sous-évaluée,
une sonde de démarrage trop courte.Les commandes
\texttt{kubectl describe} et \texttt{kubectl get events} ont été plus utiles
que toute relecture de fichier YAML, et c'est le champ \texttt{Last State}
qui a livré la clef du pod en boucle d'échecs.

La seconde tient à la mesure, en effet un déploiement conforme ne garantit pas un service
performant, et il a fallu instrumenter le système pour découvrir une forte
dégradation dès trois requêtes simultanées.Ce constat déplace la question.
Le déploiement n'est pas le facteur limitant principal dans cette expérience ;
c'est le moteur d'inférence qui devient le facteur limitant,et l'étape suivante
d'une mise en production consisterait à envisager un moteur pratiquant la
gestion paginée du cache et le batching continu.

% ============================================================
% SOURCES
% ============================================================
\section*{Sources}

\begin{itemize}
\small

\item
Kwon W. \textit{et al.},
\textit{Efficient Memory Management for Large Language Model Serving with
PagedAttention}, SOSP 2023.
\href{https://arxiv.org/abs/2309.06180}
{arxiv.org/abs/2309.06180}

\item
Yu G. \textit{et al.},
\textit{Orca, a Distributed Serving System for Transformer-Based Generative
Models}, OSDI 2022, à l'origine de l'ordonnancement par itération.

\item
Documentation vLLM,
\href{https://docs.vllm.ai}{docs.vllm.ai}, et
\textit{Inside vLLM, Anatomy of a High-Throughput LLM Inference System},
\href{https://blog.vllm.ai/2025/09/05/anatomy-of-vllm.html}
{blog.vllm.ai}

\item
Documentation Kubernetes, concepts relatifs aux Deployments, Services,
volumes persistants et sondes,
\href{https://kubernetes.io/docs/concepts/}
{kubernetes.io/docs/concepts}

\item
Documentation Ollama, API et gestion des modèles,
\href{https://github.com/ollama/ollama/blob/main/docs/api.md}
{github.com/ollama/ollama}

\item
Documentation Prometheus, types de métriques et fonction
\texttt{histogram\_quantile},
\href{https://prometheus.io/docs/}
{prometheus.io/docs}

\item
Documentation Open WebUI,
\href{https://docs.openwebui.com}
{docs.openwebui.com}

\end{itemize}

\end{document}
